\documentclass[11pt,a4paper]{article}
\usepackage[margin=23mm]{geometry}
\usepackage[T1]{fontenc}
\usepackage{lmodern,amsmath,amssymb,booktabs,longtable,array,hyperref}
\hypersetup{hidelinks}
\emergencystretch=3em
\newcommand{\A}{\textbf{(a)}\ }
\newcommand{\C}{\textbf{(c)}\ }
\newcommand{\D}{\textbf{(d)}\ }
\newcommand{\dd}{\mathrm d}
\newcommand{\file}[1]{\nolinkurl{#1}}
\newcommand{\vac}{\file{good_renorm_vacuum.py}}
\newcommand{\trc}{\file{good_renorm_trace.py}}
\newcommand{\lat}{\file{good_renorm_lattice.py}}
\newcolumntype{L}[1]{>{\raggedright\arraybackslash}p{#1}}
\title{Good Boussinesq one-loop renormalisation:\\vacuum elasticity, acoustic limits and lattice qualifications}
\author{Research note for PI Gustav Delius}
\date{4 November 2026}
\begin{document}
\maketitle
\begin{abstract}
\C Free-good-vacuum normal ordering leaves a finite curvature
$W''(0)=-9/(2\pi)$. Requiring the reference vacuum's elastic modulus to
retain its classical value fixes
$\Delta H=\hbar\frac9{4\pi}\int u^2\dd x$.
The audited, conditional equal-mode-count soliton trace then gives
$C_{\rm good}=1/(4\pi)$ and $C_{\rm good}-s_1=1/(4\pi)$, where $s_1$
is the thermodynamic acoustic coefficient. This does not complete the
projected determinant.

\C A bare-on-shell one-loop calculation has a different sequential
limit, $s_{1,\mathrm{pole,seq}}=-3/(2\pi)$ before a finite subtraction,
because soft internal phonons contribute at threshold. An exact
phonon-slope condition cannot silently replace the static condition.
\D We recommend the static reference-vacuum condition; no PI choice is
imposed. \C Consistent inverse continuation of Toda matching gives
$+3/(2\pi)$, an admissible finite local term which fails the classical
static acoustic condition. The classical Belov--Chaltikian continuum
flow has the good sign, but no canonical quantum lattice elasticity is
extracted. Fresh symbolic and high-precision checks, limitations and
failures are recorded below.
\end{abstract}

\section{Scope and the physical response being defined}
Labels mean (a) literature result, (b) literature heuristic, (c) project
derivation or check, (d) recommendation or unestablished identification.
Numerical checks are evidence, not rigorous error bounds. All new scripts
are under \file{phase2/scripts/}; they import no historical checker and
write only \file{good_renorm} outputs. Existing files and git state are
unchanged. No old scripts were run.

The controlling internal sources are \file{audit_onebody_scheme.tex},
\file{audit_det_projected.tex} and \file{errata.tex}; their qualifications
override the packages. The one-body scheme, Toda precision and projected
one-loop packages, README Status and the master report's conclusions
and ranked programme were consulted. Toda is method validation, not the
target regulator.

\paragraph{Canonical convention \C (\vac).}
Keep reference coordinates, normalization and $\hbar$ fixed:
\begin{align}
 H_g&=\int[\tfrac12\pi^2+V(u)+\tfrac12u_x^2]\dd x,
 &V(u)&=\tfrac12u^2-u^3, &u&=\phi_x,\quad\pi=\phi_t,\\
 u_t&=\partial_x\pi,
 &\pi_t&=\partial_x(u-3u^2-u_{xx}),\\
 u_{tt}&=u_{xx}-u_{xxxx}-3(u^2)_{xx}.\label{eq:pde}
\end{align}
The mean $U$ is fixed within a sector. Comparing sectors introduces a
source $J$ conjugate to $\int u\dd x$. For perturbative uniform energy
\emph{density} $e(U)$ in fixed reference length, define
\begin{equation}
 J=e'(U),\qquad K(U)=\partial J/\partial U=e''(U),\qquad
 \chi_U=\partial U/\partial J=K^{-1}.\label{eq:response}
\end{equation}
Here $\chi_U$ is strain compliance, the precise ``compressibility'' used
below. A particle-density compressibility needs an additional dictionary.
Using total energy requires dividing its second derivative by length.

\C Uniform momentum costs its canonical quadratic energy:
$\tfrac12(\pi+b)^2-\tfrac12\pi^2=b\pi+b^2/2$.
Thus the local long-wavelength theory has unit uniform inertia.
Linearizing $U_t=\pi_x$, $\pi_t=\partial_x e'(U)$ defines
$c_{\rm hyd}^2=K$. This is a thermodynamic/hydrodynamic speed.
\D Equality to the slope of an exact spectral phonon requires a specified
state and an infrared argument; Section~\ref{sec:pole} shows why the
naive on-shell perturbative replacement is insufficient.

\paragraph{The vacuum is a local perturbative branch \C (\vac).}
Putting $u=U+\eta_x$ gives quadratic elastic coefficient $m^2=1-6U$ and
\begin{equation}
 \omega(q;U)=|q|\sqrt{q^2+1-6U},\qquad U<1/6.\label{eq:disp}
\end{equation}
The sign in the question is correct. However the cubic Hamiltonian is
unbounded below even at fixed mean. On a period $2\pi$, take
$f=\cos x+\cos2x$: $\int f=0$ and $\int f^3=3\pi/2>0$.
The energy of $u=U+A f$ has a negative leading $A^3$ term, whereas its
gradient energy is quadratic. ``Vacuum'' therefore means the stable
Gaussian branch and its perturbative effective energy, not a constructed
global ground state. A stable quantum completion is extra input \D.

\section{Vacuum curvature and finite counterterms}
\paragraph{Fixed-cutoff calculation \C (\vac).}
The zero-point coefficient at a strain-independent wave-number cutoff is
\begin{align}
 W_\Lambda^{\rm bare}(U)&=\frac1{2\pi}\int_0^\Lambda
 q\sqrt{q^2+1-6U}\dd q\\
 &=\frac{(\Lambda^2+1-6U)^{3/2}-(1-6U)^{3/2}}{6\pi}.
\end{align}
Write $e=V+\hbar W+O(\hbar^2)$. The required vacuum constant and linear
counterterms are
\begin{equation}
 \Delta\mathcal H_{0,1}=-\hbar W_\Lambda^{\rm bare}(0)
 +\frac{3\hbar}{2\pi}(\sqrt{1+\Lambda^2}-1)u.\label{eq:linear}
\end{equation}
The linear term is precisely free-good-vacuum Wick ordering of $-u^3$:
\begin{equation}
 \langle u^2\rangle_0=\frac{\hbar}{2\pi}(\sqrt{1+\Lambda^2}-1),
 \qquad -:u^3:=-u^3+3\langle u^2\rangle_0u.
\end{equation}
This fixes $e(0)=e'(0)=0$ at one loop, including the finite pressure/charge
origin. The Fock reference is fixed at $U=0$, not reset for each background.

\C The finite result is
\begin{align}
 W(U)&=-\frac{(1-6U)^{3/2}-1+9U}{6\pi},\label{eq:W}\\
 W''(U)&=-\frac9{2\pi}\int_0^\infty
 \frac{q\dd q}{(q^2+1-6U)^{3/2}}
 =-\frac9{2\pi\sqrt{1-6U}}.\label{eq:Wpp}
\end{align}
The curvature is UV- and IR-convergent. Bounded, background-independent
regulator functions tending to one give the same limit by dominated
convergence. This independence holds within the specified Hamiltonian and
subtractions; it does not forbid finite local additions or changes of UV
completion. No divergent quadratic counterterm is needed for this
homogeneous one-loop calculation.

\C Independent 80-digit quadratures use the cancellation-free identity
\begin{equation}
 q\left(\sqrt{q^2+1-6U}-\sqrt{q^2+1}
       +\frac{3U}{\sqrt{q^2+1}}\right)
 =\frac{-18U^2q}{\sqrt{q^2+1}(\sqrt{q^2+1-6U}+\sqrt{q^2+1})^2}.
\end{equation}
At $U=-.08,0,.02,.10$, energy and curvature discrepancies are below
$3\times10^{-80}$. Smooth damping at $U=0$ independently verifies
\begin{equation}
 \int_0^\infty\frac{q e^{-\delta q^2}}{(1+q^2)^{3/2}}\dd q
 =1-\sqrt{\pi\delta}e^\delta\operatorname{erfc}\sqrt\delta\longrightarrow1.
\end{equation}
Finite regulator error is distinct from quadrature error.

\paragraph{Classical reference-vacuum elasticity \C (\vac).}
Adding $\Delta\mathcal H=\hbar\lambda u^2$ to the normal-ordered reference
produces
\begin{align}
 K(U)&=1-6U+\hbar\left(2\lambda-
       \frac9{2\pi\sqrt{1-6U}}\right)+O(\hbar^2),\\
 \chi_U&=\frac1{1-6U}-\hbar
 \frac{2\lambda-9/(2\pi\sqrt{1-6U})}{(1-6U)^2}+O(\hbar^2),\\
 c_{\rm hyd}(0)&=1+\hbar s_1+O(\hbar^2),
 &s_1&=\lambda-\frac9{4\pi}.\label{eq:s1}
\end{align}
Consequently $K(0)=1+O(\hbar^2)$, $\chi_U(0)=1+O(\hbar^2)$ and
$c_{\rm hyd}(0)=1+O(\hbar^2)$ are equivalent and fix
\begin{equation}
 \boxed{\lambda_{\rm acoustic}=\frac9{4\pi}
 =0.71619724391352901096\ldots .}\label{eq:lambda}
\end{equation}
Normal ordering alone instead prescribes $\lambda=0$, leaving
$s_1=-9/(4\pi)$. The finite bubble is not a Wick tadpole.
These expansions are not uniform near $U=1/6$.

\paragraph{Preserving the entire background family \C (\vac).}
A constant $\lambda$ cannot cancel the $U$ dependence in \eqref{eq:Wpp}.
Requiring $K(U)=1-6U+O(\hbar^2)$ throughout $U<1/6$, with the same
energy/pressure origins, instead requires
\begin{equation}
 \Delta\mathcal H=\hbar G(u),\quad
 G(u)=-W(u)=\frac{(1-6u)^{3/2}-1+9u}{6\pi}
 =\frac9{4\pi}u^2+\frac9{4\pi}u^3+\frac{81}{16\pi}u^4+\cdots .
 \label{eq:G}
\end{equation}
This is an additional finite nonlinear condition, not a UV necessity.
The closed function is real only for $u<1/6$; its small-field expansion
specifies local vertices, not a completion for arbitrary cores or fields.
Its quadratic coefficient suffices for the small-amplitude table; using
only $\lambda u^2$ does not implement it at finite amplitude.

\section{Exact phonon slope versus a sequential perturbative pole}
\label{sec:pole}
\paragraph{Status of the candidate \D.}
In a specified stable quantum completion one could impose the classical
slope of the low-$q$ spectral edge or phonon pole. This cubic model supplies
neither that completion nor a proof of an isolated pole. If hydrodynamics
identifies its slope with $\sqrt K$, the condition gives \eqref{eq:lambda}.
That is an explicit assumption. Evaluating the one-loop self-energy on
the bare dispersion before sending $q$ to zero is a different limit.

\paragraph{A direct infrared test \C (\vac).}
Expansion of $\frac12\operatorname{Tr}\log(D_0+\delta D)$ to second order
in a background insertion gives
\begin{equation}
 \frac{\Sigma_E(\Omega,q)}{\hbar q^2}
 =-18\int\frac{\dd p\dd\nu}{(2\pi)^2}
 \frac{p^2(q-p)^2}{(\nu^2+\omega_p^2)
 ((\Omega-\nu)^2+\omega_{q-p}^2)}.\label{eq:bubble}
\end{equation}
Here each insertion is $-6u$ times the two spatial derivatives. The factor
$-18$ is the corresponding quadratic-kernel coefficient.
The frequency convolution, independently integrated numerically, is
\begin{equation}
 \int\frac{\dd\nu}{2\pi}
 \frac1{(\nu^2+a^2)((\Omega-\nu)^2+b^2)}
 =\frac{a+b}{2ab((a+b)^2+\Omega^2)}.
\end{equation}
Its static limit reproduces \eqref{eq:Wpp}. Continuing
$\Omega=-i\omega_q$ leaves the integral without the $-18$ factor,
\begin{equation}
 I(q)=\int_{\mathbb R}\frac{\dd p}{2\pi}
 \frac{p^2(q-p)^2(\omega_p+\omega_{q-p})}
 {2\omega_p\omega_{q-p}[(\omega_p+\omega_{q-p})^2-\omega_q^2]}.
\end{equation}
For $q>0$, a split at $p=0,q$ gives
\begin{equation}
 \lim_{q\downarrow0}I_{0<p<q}=-\frac1{12\pi},\qquad
 \lim_{q\downarrow0}I_{p<0\,\text{or}\,p>q}=\frac1{4\pi},\qquad
 \lim I(q)=\frac1{6\pi}.\label{eq:soft}
\end{equation}
Putting $p=qx$ makes the interior integrand, including $\dd p$ but before
$1/(2\pi)$, tend to $-\dd x/6$. The exterior tends to the static bubble.
The endpoint numerator zeros make this integral finite. The script
rationalizes the denominator to avoid cancellation and checks the soft
scaling symbolically as well as by split quadrature.

\C The resulting sequential pole slope is
\begin{equation}
 s_{1,\mathrm{pole,seq}}=-9\lim I(q)+\lambda
 =-\frac3{2\pi}+\lambda,\label{eq:pole}
\end{equation}
which is not $s_1$ in \eqref{eq:s1}. At $\lambda=0$, the finite-$q$
value is $-0.477443610314\ldots$ at $q=.01$; its residual from the limit
is below $3\times10^{-25}$ at $q=10^{-12}$ (asserted at $10^{-21}$).
The order of limits is expansion in $\hbar$ at fixed $q$, then $q\to0$.
The bare collinear frequency splitting is $O(q^3)$, while a one-loop
velocity shift changes frequency by $O(\hbar q)$; this threshold expansion
loses uniformity when $q^2$ is comparable to $\hbar$.
\D Resummation or an exact low-energy solution is needed for the strict
fixed-$\hbar$, $q\to0$ spectral slope. No value of $\lambda$ is asserted
to settle that exact condition unconditionally.
\C Formally cancelling \eqref{eq:pole} chooses $3/(2\pi)$ and leaves a
nonclassical static elasticity. The table retains this diagnostic.

\section{Conditional one-body result and candidate table}
\paragraph{Trace and fixed momentum \C (\trc).}
Retain the audited equal-mode-count endpoint and finite Wick pressure
origin. For the stable small-amplitude branch $0<k<\sqrt3/2$,
\begin{align}
 F_0(k)&=\frac{\sqrt3}{\pi}(\tfrac34-k^2)
 \operatorname{artanh}\frac{2k}{\sqrt3}-\frac{3k}{2\pi},\\
 P_0&=\epsilon\frac23k^3\sqrt{1-k^2},\qquad
 E_0=\frac2{15}k^3(5-4k^2),\qquad \dd E_0=v\dd P_0.
\end{align}
Here $\hbar F_0$ is the conditional energy correction at fixed momentum.
Fresh convergent quadratures of the energy and momentum phase densities
reproduce it at four independent widths to below $10^{-65}$. The endpoint
and tadpole subtraction are combined in the convergent energy integrand;
no fitted target sets the constant.

\C The classical moments are $\int u^2=2k^3/3$,
$\int u^3=4k^5/15$ and $\int u_x^2=2k^5/15$.
At fixed $P$, stationarity of $H_0-vP$ removes first-order profile
variation. Hence
\begin{align}
 F_\lambda(k)&=F_0(k)+\frac{2\lambda}{3}k^3,\label{eq:Flambda}\\
 F_0(k)&=-\frac4{3\pi}k^3-\frac{16}{45\pi}k^5+O(k^7),\\
 C_{\rm good}(\lambda)&=\lim_{k\to0}\frac{F_\lambda(k)}{|P_0(k)|}
 =\lambda-\frac2\pi,\qquad
 \boxed{C_{\rm good}-s_1=\frac1{4\pi}.}\label{eq:C}
\end{align}
Independent differentiation of the exact classical soliton with stiffness
$a=1+2\hbar\lambda$ at fixed $P$ verifies this shift.
A finite linear term $\hbar\mu\int u$ instead produces $3\mu/k^2$ in the
ratio and destroys its finite limit. This fixes the need for a separate
pressure/charge-origin convention. The higher terms $u^3,u_x^2$ contribute
at $k^5$ and do not change $C$.

\begin{center}\small
\begin{tabular}{@{}L{5.3cm}cccc@{}}\toprule
Condition / prescription & $\lambda$ & $s_1$ & $C_{\rm good}$ & $C_{\rm good}-s_1$\\\midrule
Free good-vacuum normal ordering & $0$ & $-9/(4\pi)$ & $-2/\pi$ & $1/(4\pi)$\\[3pt]
Classical $K(0)$ or $\chi_U(0)$ & $9/(4\pi)$ & $0$ & $1/(4\pi)$ & $1/(4\pi)$\\[3pt]
Classical hydrodynamic speed at $U=0$ & $9/(4\pi)$ & $0$ & $1/(4\pi)$ & $1/(4\pi)$\\[3pt]
Classical $K(U)$ for all backgrounds; full $G(u)$ required & $9/(4\pi)$\footnotemark & $0$ & $1/(4\pi)$ & $1/(4\pi)$\\[3pt]
Inverse continuation of Toda matching & $3/(2\pi)$ & $-3/(4\pi)$ & $-1/(2\pi)$ & $1/(4\pi)$\\[3pt]
Copy bad coefficient without continuing density & $-3/(2\pi)$ & $-15/(4\pi)$ & $-7/(2\pi)$ & $1/(4\pi)$\\[3pt]
Cancel sequential bare-on-shell slope only & $3/(2\pi)$ & $-3/(4\pi)$ & $-1/(2\pi)$ & $1/(4\pi)$\\[3pt]
Exact spectral slope without infrared identification & unknown & --- & --- & ---\\[3pt]
Canonical good quantum lattice matching & unknown & --- & --- & ---\\\bottomrule
\end{tabular}
\end{center}
\footnotetext{Quadratic coefficient of $G$, not its complete specification.}
\C All formula entries are checked in \vac\ and \trc; all $C$ values
are conditional on the audited trace. The exact spectral row becomes the
hydrodynamic row only under the stated identification.
\D Recommendation: use $K(0)=1$ with $e(0)=e'(0)=0$ for the next
continuum comparison. This condition is reproducible and defined by the
stable reference vacuum. It is not forced by integrability or Toda and is
not recorded as the PI's adopted scheme.

\paragraph{Machine-readable handoff.}
\C \file{scripts/good_renorm_conditions.json} contains exact strings and
70-digit decimal $\lambda$ values, all table rows, the reference
subtractions, $F_\lambda$, $C(\lambda)$ and the definition of $s_1$.
The recommendation is labelled (d), \file{adopted_condition} is null,
and \file{physical_determinant_completed} is false. The two-body package
can use $\lambda$ parametrically or quote the recommended conditional value.
The invariant is restricted to common quadratic bulk subtractions with
units, kinetic normalization, pressure origin and boundary prescription
fixed. Invariance under all counterterms or quantum completions is not
claimed. The full $G$ prescription needs its extra vertices at finite
amplitude, not just the exported quadratic coefficient.

\section{Continuation test: admissible term, failed acoustic match}
\paragraph{Consistent continuation \C (\trc).}
The audited branch uses $k_g=ik_b$, $\alpha=i\beta$, $x_g=-ix_b$,
$U_g=-U_b$ and $F_b(k)=-iF_g(ik)$.
The velocities continue by $\cos(i\beta)=\cosh\beta$; the stipulated
spectral conventions have $\ell_b=-\ell_g(i\beta)$.
The density rule, not rapidity substitution alone, determines
\begin{equation}
 W_b(U)=-W_g(-U),\qquad \lambda_b=-\lambda_g.
\end{equation}
Therefore $\lambda_b=-3/(2\pi)$ inversely continues to
$\lambda_g=+3/(2\pi)$, not the same signed number.
Continuing the Toda vacuum function gives
\begin{equation}
 W_{T\to g}(U)=-\frac{e^{-3U}-1+3U}{6\pi},\qquad
 \tfrac12W_{T\to g}''(0)=-\frac3{4\pi}.
\end{equation}
The independently recomputed unmodified good vacuum has
$s_{1,g}=-9/(4\pi)$; their difference is $3/(2\pi)$.
Inverse-Toda matching therefore leaves nonzero thermodynamic sound shift
$-3/(4\pi)$ and conditional $C=-1/(2\pi)$. It fails the classical good
acoustic condition by $-3/(4\pi)$ in both columns.

\C This real local finite term has no algebraic prohibition; for
sufficiently small $\hbar$ its vacuum stiffness remains positive. Thus
``cannot define any legitimate perturbative good scheme'' is too strong.
What fails is identifying it with classical-acoustic renormalisation or
deriving it from a real good-equation lattice.
\D Physical selection or a stronger no-go needs a quantum completion.
The coincidence with the sequential-pole subtraction is only an algebraic
equality, not a proof of an exact-slope or lattice interpretation.

\paragraph{No real good analogue of the bad turning-band mechanism
\C (\trc).}
For $q>0$,
\begin{equation}
 \omega_g'(q)=\frac{1+2q^2}{\sqrt{1+q^2}},\qquad
 (\omega_g')^2-1=\frac{q^2(3+4q^2)}{1+q^2}>0.
\end{equation}
There is no group-velocity match with a subsonic right-moving soliton;
negative $q$ has the opposite group velocity. Thus the bad-band mechanism
does not transfer to the real good spectrum. The controlling audit also
warns that light-root logarithm cancellation does not prove that an entire
trace is supported only near a short-wave root. Neither observation
forbids a finite $u^2$ scheme change in the good theory.

\section{Integrable lattice assessment}
\paragraph{Reading scope and source facts.}
\C Both abstract pages were opened. Belov--Chaltikian's HTML full-text
sections 1--6 were read, especially sections 4--5 \cite{BC}.
For Hikami \cite{Hikami}, PDF passages were read in section 1,
section 2.1 (commutators and quasi-periodic condition), and the openings
of sections 5 and 7. Its full 25-page proof was not read. Its experimental
HTML failed to load, but the PDF was accessible. Other discretisation
papers returned by a brief search are not used as evidence here.

\A The sources exhibit a lattice Boussinesq hierarchy and a classical
continuum flow. Hikami constructs a Baxter operator for the quantum
chain and uses chiral-field commutators in the continuum discussion.
These algebraic results do not identify the canonical vacuum elasticity
needed here. Belov--Chaltikian's displayed time convention gives
\begin{equation}
 U_T=-U_{XX}+2W_X,\qquad
 W_T=W_{XX}-\tfrac23U_{XXX}-\tfrac23UU_X.\label{eq:source}
\end{equation}
Hikami reverses both right-hand sides.

\paragraph{A real good-sign PDE dictionary \C (\lat).}
Eliminating $W$ gives
\begin{equation}
 U_{TT}=-\tfrac13U_{XXXX}-\tfrac23(U^2)_{XX}.
\end{equation}
The real transformation
\begin{equation}
 X=x,\quad T=\sqrt3\,t,\quad U=-\tfrac14+\tfrac32u,
 \quad \pi=\frac{4\sqrt3}{3}(W-U_X/2)\label{eq:latticedict}
\end{equation}
produces \eqref{eq:pde} and $u_t=\pi_x$. Reverse $T$ for Hikami's sign.
Thus the good fourth-derivative sign is available at the level of the
printed classical continuum flow. This is a PDE dictionary, not a
symplectic/action or $\hbar$ dictionary.

\paragraph{Linear lattice diagnostic and a failed shortcut \C (\lat).}
The base flow printed in both sources is
\begin{equation}
 \dot L_n=W_{n-1}-W_n+L_n(L_{n+1}-L_{n-1}),\qquad
 \dot W_n=W_n(L_{n+2}-L_{n-1}).
\end{equation}
Linearization around $(L_0,W_0)$ with $z=e^{iK}$ gives
\begin{equation}
 A(z)=\begin{pmatrix}L_0(z-z^{-1})&z^{-1}-1\\
 W_0(z^2-z^{-1})&0\end{pmatrix},\quad
 D=-4L_0^2\sin^2K+8W_0(\cos K-\cos2K).
\end{equation}
At $(1/3,1/27)$,
\begin{equation}
 D=-\frac4{27}(1-\cos K)^2,\qquad
 r_\pm=\frac{i\sin K}{3}\pm\frac{i(1-\cos K)}{3\sqrt3}.
\end{equation}
This base-flow critical background has oscillatory linear modes. That
alone is not a positive-energy quantum vacuum. Inserting the printed
leading background expansions $L_0=(1-\Delta^2U/3)/3$,
$W_0=(1-\Delta^2U)/27$, $K=\Delta q$ gives
\begin{equation}
 D=-\frac{\Delta^4}{27}(q^4+4Uq^2)+O(\Delta^6).
\end{equation}
Its acoustic $U$ sign disagrees with a naive identification of these
\emph{base-flow} frequencies with the shifted target in
\eqref{eq:latticedict}. This refutes that shortcut, not the higher-flow
construction: Hikami explicitly specifies a higher flow, and
Belov--Chaltikian takes regularized Hamiltonian combinations before the
continuum limit.

\paragraph{Missing ingredients for $\lambda$ \D.}
A matching calculation still needs:
\begin{enumerate}
\item A self-adjoint real form and the precise commuting-Hamiltonian
combination generating the desired higher flow.
\item Its stable uniform state and physical positive-frequency modes over
the entire Brillouin zone, as functions of the mean strain.
\item The symplectic/action conversion to
$\Omega=\int\dd\pi\wedge\dd\phi$, fixing $\hbar$ and the kinetic inertia.
\item The complete lattice vacuum energy with fixed mode count, converted
to reference-length density, and its constant/pressure subtractions.
\end{enumerate}
A PDE match and a Baxter operator alone do not give those identifications.
\C Once a canonical lattice $s_{1,L}$ is supplied, \eqref{eq:s1} fixes
$\lambda=s_{1,L}+9/(4\pi)$. No such numerical $s_{1,L}$ or $\lambda$ is
extracted here. \D The classical good-sign result makes this a plausible
route, not a completed quantum regulator or a no-go theorem.

\section{Checks, error budget and failures recorded}
\small
\begin{longtable}{@{}L{4.3cm}L{6.3cm}L{3.5cm}@{}}\toprule
Claim & Evidence and qualification & Checker\\\midrule\endhead
\C Vacuum sign, curvature, compliance, counterterms & SymPy identities;
80-digit subtracted and smooth-regulator quadratures at tolerance $10^{-70}$ & \vac\\[4pt]
\C Nonuniform on-shell limit & Frequency convolution, symbolic soft
scaling, split integral down to $q=10^{-12}$; limit tolerance $10^{-21}$ & \vac\\[4pt]
\C Conditional $F_\lambda$, $C$, invariant & Exact moments, series and
fixed-$P$ deformation; independent trace quadratures at tolerance $10^{-65}$ & \trc\\[4pt]
\C Continuation sign and mismatch & Density continuation and vacuum
curvatures; group-velocity identity & \trc\\[4pt]
\C Lattice good sign and base-flow limitation & PDE elimination,
field/time map, lattice linearization and 80-digit eigenvalues & \lat\\[4pt]
\D Physical determinant, exact spectral slope, lattice matching & Not
proved; unknown physical identifications have no numerical error bar & Open\\\bottomrule
\end{longtable}
\normalsize

\paragraph{Error budget \C.}
The recommended conditional numbers are exact expressions,
$\lambda=9/(4\pi)$ and $C=1/(4\pi)=0.07957747154594766788\ldots$,
not fits. Numerical code uses 80 decimal digits. Vacuum quadrature errors
are below $3\times10^{-80}$, trace discrepancies below $10^{-80}$
(asserted at $10^{-65}$). Four-width interpolation at $k=10^{-5}/2^j$
checks the $-2/\pi$ reference limit to better than $10^{-40}$.
The sequential-pole finite-$q$ residual is below $3\times10^{-25}$ at
$q=10^{-12}$, separate from arithmetic accuracy. These are deterministic
diagnostics, not rigorous rounding bounds. Unknown determinant boundary
contributions, quantum UV matching and nonuniform infrared or
finite-$\hbar$ soliton limits are structural qualifications. No
higher-loop error estimate is supplied.

\paragraph{Failures recorded.}
\begin{enumerate}
\item \C Exploratory direct on-shell quadrature raised
\file{ZeroDivisionError} at $q=10^{-4}$ from subtracting nearly equal
squared frequencies. Exact denominator rationalization repaired the
calculation, with no weakened tolerance or dropped integration region.
\item \C Identifying the sequential bare-on-shell slope with static
elasticity fails by $3/(4\pi)$. The discrepancy is retained, not removed
by relabelling $s_1$.
\item \C Copying Toda's signed term fails consistent density continuation
and classical good-acoustic matching. Correct inverse continuation still
fails the latter; it is not ruled out as an arbitrary finite local term.
\item \C The base-lattice-flow/background shortcut fails the acoustic
sign comparison. No vacuum elasticity is inferred from it.
\item \C Hikami's HTML fetch failed; only the listed PDF passages and
abstract are claimed as read.
\item \C The first report-writing command used a redundant
\file{phase2/} prefix while already in that directory. The file was not
created and the attempted build could not find it. The temporary
\file{texput.log} was removed and the path corrected.
\item \C A PDF-render attempt started before asynchronous builds had
finished and reported a missing trailer/xref. Both processes were awaited;
the final two builds were rerun strictly sequentially before rendering.
\end{enumerate}
No assertion in the final three scripts was weakened. Each prints
\texttt{ALL CHECKS PASSED} after its assertions. This certifies only those
assertions, not unresolved physical identifications. The report is built
twice with \file{pdflatex} in \file{phase2/}; rendered pages are inspected.

\begin{thebibliography}{9}
\bibitem{BC} A.~A.~Belov and K.~D.~Chaltikian,
\emph{Lattice analogues of W-algebras and Classical Integrable Equations},
\emph{Phys. Lett. B} \textbf{309} (1993) 268--274,
\href{https://arxiv.org/abs/hep-th/9303166}{arXiv:hep-th/9303166}.
Abstract and HTML full-text sections 1--6 read; equations used from
sections 4--5.
\bibitem{Hikami} K.~Hikami,
\emph{Baxter Equation for Quantum Discrete Boussinesq Equation},
\emph{Nucl. Phys. B} \textbf{604} (2001) 580--602,
\href{https://arxiv.org/abs/nlin/0102021}{arXiv:nlin/0102021}.
Abstract and specified PDF excerpts read; not the complete text.
\end{thebibliography}
\end{document}
